% ---------------------------------------------------------------
% Two-page extended abstract
% Excess Separability: Nuisance-Controlled Residual-Stream Probing
% for Benchmark Contamination Detection
%
% Compile with: pdflatex abstract_of_thesis.tex (run twice)
% Self-contained: no .bib file, no ACL style file required.
% ---------------------------------------------------------------
\documentclass[10pt,letterpaper]{article}

\usepackage[letterpaper,top=0.6in,bottom=0.6in,left=0.7in,right=0.7in]{geometry}
\usepackage[T1]{fontenc}
\usepackage[utf8]{inputenc}
\usepackage{times}
\usepackage{amsmath}
\usepackage{amssymb}
\usepackage{microtype}
\usepackage{multicol}
% Measured values, generated by render_all.py.
% Auto-generated by render_all.py. Every value is measured.
\newcommand{\PropAGrowth}{1.00}
\newcommand{\PropACorr}{0.01}
\newcommand{\PropAKurt}{3.0}
\newcommand{\PropAEig}{0.025}
\newcommand{\PropBGrowth}{1.75}
\newcommand{\PropBCorr}{0.65}
\newcommand{\PropBKurt}{9.3}
\newcommand{\PropBEig}{0.188}
\newcommand{\CapAContrast}{(96,0.02) (400,0) (1200,0.025)}
\newcommand{\CapAMaxLevel}{(96,0.03) (400,0.53) (1200,0.9875)}
\newcommand{\CapAMeanLevel}{(96,0.03) (400,0.64) (1200,0.975)}
\newcommand{\CapANaive}{(96,1) (400,1) (1200,1)}
\newcommand{\CapContrastMatch}{0.020}
\newcommand{\CapMaxMatch}{0.030}
\newcommand{\CapNaiveMatch}{1.000}
\newcommand{\CapBAnMatch}{0.652}
\newcommand{\CapContrastLarge}{0.000}
\newcommand{\CapMaxLarge}{0.530}
\newcommand{\CapNaiveLarge}{1.000}
\newcommand{\CapBAnLarge}{0.609}
\newcommand{\CapContrastHuge}{0.025}
\newcommand{\CapMaxHuge}{0.988}
\newcommand{\CapNaiveHuge}{1.000}
\newcommand{\CapBAnHuge}{0.570}
\newcommand{\BootMin}{0.030}
\newcommand{\BootMax}{0.087}
\newcommand{\CapBootMatch}{0.080}
\newcommand{\CapBootLarge}{0.030}
\newcommand{\CapBootHuge}{0.087}
\newcommand{\CapPdimLo}{96}
\newcommand{\CapPdimHi}{1200}
\newcommand{\CapRatio}{12}
\newcommand{\SlopeARaw}{(0,0.01667) (0.25,0.2833) (0.5,0.7167)}
\newcommand{\SlopeAAdj}{(0,0.08333) (0.25,0.05) (0.5,0.05)}
\newcommand{\SlopeARawZero}{0.017}
\newcommand{\SlopeAAdjZero}{0.083}
\newcommand{\SlopeARawTwentyFive}{0.283}
\newcommand{\SlopeAAdjTwentyFive}{0.050}
\newcommand{\SlopeARawFifty}{0.717}
\newcommand{\SlopeAAdjFifty}{0.050}
\newcommand{\SlopeBRaw}{(0,0) (0.25,0) (0.5,0)}
\newcommand{\SlopeBAdj}{(0,0.08333) (0.25,0.08333) (0.5,0.03333)}
\newcommand{\SlopeBRawZero}{0.000}
\newcommand{\SlopeBAdjZero}{0.083}
\newcommand{\SlopeBRawTwentyFive}{0.000}
\newcommand{\SlopeBAdjTwentyFive}{0.083}
\newcommand{\SlopeBRawFifty}{0.000}
\newcommand{\SlopeBAdjFifty}{0.033}
\newcommand{\PowerA}{(0,0.02) (0.5,0.04) (1,0.38) (1.5,0.86) (2,0.975) (3,1)}
\newcommand{\PwAZero}{0.020}
\newcommand{\PwAFive}{0.040}
\newcommand{\PwATen}{0.380}
\newcommand{\PwAFifteen}{0.860}
\newcommand{\PwATwenty}{0.975}
\newcommand{\PwAThirty}{1.000}
\newcommand{\PowerB}{(0,0.06) (0.5,0.14) (1,0.34) (1.5,0.78) (2,0.975) (3,1)}
\newcommand{\PwBZero}{0.060}
\newcommand{\PwBFive}{0.140}
\newcommand{\PwBTen}{0.340}
\newcommand{\PwBFifteen}{0.780}
\newcommand{\PwBTwenty}{0.975}
\newcommand{\PwBThirty}{1.000}
\newcommand{\ProfNullA}{(0,0.7028) (1,0.708) (2,0.71) (3,0.7065) (4,0.7078) (5,0.7105) (6,0.7012) (7,0.7105) (8,0.7105) (9,0.7085) (10,0.7107) (11,0.7007) (12,0.7017)}
\newcommand{\ProfPlcA}{(0,0.6782) (1,0.6666) (2,0.6701) (3,0.6756) (4,0.6682) (5,0.6739) (6,0.661) (7,0.669) (8,0.6733) (9,0.6726) (10,0.6752) (11,0.6629) (12,0.6657)}
\newcommand{\ProfSigA}{(0,0.7028) (1,0.7077) (2,0.7172) (3,0.7173) (4,0.729) (5,0.7543) (6,0.7807) (7,0.7948) (8,0.8077) (9,0.7998) (10,0.7795) (11,0.7464) (12,0.7285)}
\newcommand{\ProfNullFirstA}{0.703}
\newcommand{\ProfNullLastA}{0.702}
\newcommand{\ProfArgmaxA}{8}
\newcommand{\ProfNullB}{(0,0.7539) (1,0.7473) (2,0.7406) (3,0.7318) (4,0.7222) (5,0.716) (6,0.7113) (7,0.7001) (8,0.692) (9,0.6837) (10,0.6762) (11,0.6678) (12,0.6572)}
\newcommand{\ProfPlcB}{(0,0.7306) (1,0.7253) (2,0.7146) (3,0.7054) (4,0.6961) (5,0.684) (6,0.6785) (7,0.6676) (8,0.6579) (9,0.6499) (10,0.6428) (11,0.6317) (12,0.6239)}
\newcommand{\ProfSigB}{(0,0.7492) (1,0.7443) (2,0.7337) (3,0.7237) (4,0.7237) (5,0.7297) (6,0.7461) (7,0.7556) (8,0.7499) (9,0.739) (10,0.7311) (11,0.7204) (12,0.7142)}
\newcommand{\ProfNullFirstB}{0.754}
\newcommand{\ProfNullLastB}{0.657}
\newcommand{\ProfArgmaxB}{7}
\newcommand{\SizeMatchedA}{0.067}
\newcommand{\SizeHalfA}{0.217}
\newcommand{\SizeMatchedB}{0.040}
\newcommand{\SizeHalfB}{0.040}
\newcommand{\KeyFPR}{(0.5,0.05) (2,0.025) (5,0.025)}
\newcommand{\KeyPower}{(0.5,0.825) (2,0.65) (5,0.45)}
\newcommand{\KeyFprFive}{0.050}
\newcommand{\KeyPwFive}{0.825}
\newcommand{\KeyFprTwenty}{0.025}
\newcommand{\KeyPwTwenty}{0.650}
\newcommand{\KeyFprFifty}{0.025}
\newcommand{\KeyPwFifty}{0.450}
\newcommand{\TimeSetup}{6.51}
\newcommand{\TimePerPerm}{1.19}
\newcommand{\TimeB}{2.4}
\newcommand{\TimeRefitH}{1.9}
\newcommand{\TimeSpeedup}{2919}
\newcommand{\TimeN}{1000}
\newcommand{\TimeD}{2048}
\newcommand{\RunMinutes}{10}
\newcommand{\NPerSet}{200}
\newcommand{\BPerm}{300}

\usepackage[numbers]{natbib}
\usepackage[hidelinks]{hyperref}

\setlength{\parskip}{0pt}
\setlength{\parindent}{1em}
\setlength{\abovedisplayskip}{3pt plus 1pt minus 1pt}
\setlength{\belowdisplayskip}{3pt plus 1pt minus 1pt}
\makeatletter
\renewcommand\section{\@startsection{section}{1}{\z@}%
  {-1.4ex \@plus -0.4ex \@minus -.2ex}{0.5ex \@plus .1ex}%
  {\normalfont\normalsize\bfseries}}
\makeatother

\newcommand{\BA}{\mathrm{BA}}
\newcommand{\nuis}{\mathrm{nuis}}

\pagestyle{empty}

\begin{document}

\begin{center}
{\large\bfseries Excess Separability: Nuisance-Controlled Residual-Stream Probing\\[1pt]
for Benchmark Contamination Detection}\\[4pt]
{\normalsize Florian Braun}\quad
{\small\texttt{\href{mailto:braunf25@proton.me}{braunf25@proton.me}}}\quad
{\small\texttt{\href{https://braunf.com/}{braunf.com}}}\\[2pt]
{\small Extended abstract of a methodological proposal paper}\\[3pt]
{\footnotesize Code and artefacts:
\href{https://github.com/mabushi-lab/residual-stream-contamination-probing}%
{github.com/mabushi-lab/residual-stream-contamination-probing}\quad
\href{https://doi.org/10.5281/zenodo.21969647}{doi:10.5281/zenodo.21969647}}
\end{center}

\vspace{0pt}

\section{Background and Motivation}

A benchmark score is a claim about capability, and contamination breaks the claim by making the score partly a report of what a model has memorised. The three established detection routes each need something usually missing. Overlap search needs the training corpus, which for most released models is not public. Likelihood-based membership inference needs only the model but fixes a scalar test statistic in advance, and under exchangeable splits performs close to chance at pretraining scale \cite{duan2024mia}. Canary insertion works but must be planned before publication, which does nothing for benchmarks already in use.

A fourth route has appeared. \citet{liu2024probing} train a linear probe on internal activations and report state-of-the-art AUC on WikiMIA and an arXiv-derived set. The motivation is sound: a probe learns the discriminative direction instead of assuming which output statistic carries it. The problem is the reported quantity. A probe separating two item sets will use anything that separates them, and \citet{das2024blind} show that on the standard evaluation sets a classifier with no access to the target model achieves high AUC, because member and non-member sets were split by publication date and so differ in topic and style. Reported accuracy is thus the sum of memorisation and distribution shift, with the split guaranteeing the second term is large. The calibration point sits in \citeauthor{liu2024probing}'s own paper: on the randomised-split contaminated checkpoints of \citet{oren2024proving} their probe reaches 54.0 and 51.9 AUC at duplication counts 1 and 2, against a chance baseline of 50.0. The instrument is plausible; the protocol around it is not.

\section{Research Questions}

\textbf{RQ1.} How much of the separability reported by activation-based detectors is surface difference rather than memorisation? \textbf{RQ2.} What is the correct statistic for a probe-based contamination claim, and what null does it need? \textbf{RQ3.} Does the extraction point, before or after the answer tokens, change what the detector estimates? \textbf{RQ4.} At what duplication count does a familiarity signal become detectable?

\section{Methodology}

The proposal, Residual-Stream Contamination Probing (RSCP), is a set-level two-sample test rather than an item-level classifier, in the tradition of dataset inference. Five components.

\textit{Exchangeability requirement.} A suspect set $S$ and reference set $R$ are admissible only if, absent contamination, an analyst shown item text alone could not beat chance at saying which set it came from. Three constructions qualify: randomised injection before training (exact); a commissioned twin such as GSM1k against GSM8K \cite{zhang2024gsm1k}; a within-corpus held-out split such as Pile train against validation for Pythia \cite{biderman2023pythia}, known to fail on drifting subdomains. Temporal and cross-source splits, covering most existing evaluation sets, do not qualify.

\textit{Prefix extraction rule.} Activations are read at the item prefix, upstream of any answer tokens. A state computed over the full item is a sufficient statistic for the model's scoring of the gold answer, so a probe on it can approximate whether the model answers correctly, and the detector then measures the inflation it was meant to explain. This is the structure of the pre- versus post-continuation artefact identified in the author's earlier layer-wise work on Japanese-English pragmatic reasoning \cite{braun2026crosslingual}.

\textit{A declared nuisance family.} Separability is reported against an explicit control: mean-pooled layer-0 embeddings, character $n$-gram TF-IDF, surface statistics, and likelihood summaries under an independent reference model. That last block deliberately includes the statistics likelihood-based membership inference uses, computed on a model that cannot have memorised the items.

\textit{A depth-profile contrast, not a level.} The obvious statistic is
excess separability $\Delta_\ell = \BA_\ell - \BA_{\nuis}$. We measured it and
discarded it: at matched degrees of freedom a smaller control block estimates
its direction more efficiently, so under a true null the level statistic
rejects \CapMaxMatch{} of the time against a dimension-matched control and
\CapMaxHuge{} against one of $p=\CapPdimHi$. A zero-sum contrast
$T_c=\sum_\ell w_\ell \BA_\ell$ cancels the control term algebraically and
tests the shape of the profile instead of its level.

\textit{A level-matched placebo baseline.} Contrasting against a flat profile
also fails, in both directions: with surface decodability rising across depth
the uncorrected contrast rejects a true null \SlopeARawFifty{} of the time, and
on a residual stream where the surface signal is written once while noise
accumulates the profile \emph{declines} (\ProfNullFirstB{} to
\ProfNullLastB{}) and the test loses all power. The placebo split already
estimates that baseline, since both its halves are non-members. Recentring on
it, $T_{\mathrm{adj}}=\sum_\ell w_\ell(\BA_\ell-\BA^{\mathrm{plc}}_\ell)$,
replaces an untestable global assumption with a per-model measurement. The
baseline must match the observed comparison in sample size, which is why the
protocol requires $|R|=2|S|$, and in embedding-layer separability, which is
found by search.

\textit{A permutation null that refits the probe.} The probe is cross-fitted
ridge regression, chosen because out-of-fold fitted values are linear in the
labels, $\hat u_\ell = \Pi_\ell\tilde y$ with $\Pi_\ell$ label-free. An item
bootstrap holds the fit fixed and rejects \BootMin{} to \BootMax{} of true
nulls;
permuting labels against $\Pi_\ell$ refits the direction and is calibrated,
at one matrix-vector product per draw, a factor of \TimeSpeedup{} against
refitting.

\textit{Exposure calibration.} On models trained with known duplication counts $m$, fit $g(m)=\gamma(1-e^{-\beta m})$ and invert, giving an exposure interval and a detection floor $m^{*}$ rather than a verdict.

\section{Results and Worked Examples}

\textit{Measured.} The statistical machinery was validated on synthetic
activations under two deliberately dissimilar generative models: independent
Gaussian layers, and a residual stream with anisotropic power-law covariance,
heavy-tailed innovations, growing norm ($\times\PropBGrowth$) and strongly
correlated layers (\PropBCorr{} against \PropACorr). The study falsified three
of our own design choices, each described above. The corrected protocol holds
its nominal error rate throughout: \SlopeAAdjFifty{} where the uncorrected
contrast reaches \SlopeARawFifty{}, and \CapContrastMatch{} to
\CapContrastHuge{} across a \CapRatio-fold change in control-set dimension. A
half-size baseline gives \SizeHalfA{} against \SizeMatchedA{} matched. Power
reaches \PwAFifteen{} and \PwBFifteen{} at $\varepsilon=1.5$ under the two
simulators and saturates by $\varepsilon=3$. Degrading the analyst's surface
key costs power (\KeyPwFive{} to \KeyPwFifty{}) and not validity, so the
failure mode is blindness rather than error. The whole study is
\RunMinutes{} minutes on four CPU cores.

\textit{Worked examples.} These are proposals. Numbers cited below are published and attributed; the paper's own tables carry synthetic values marked illustrative, present only to fix the reporting format.

\textit{Pythia and the Pile.} Pile train against validation, per subdomain. The exact data order is public, so items can be stratified by the step at which they were seen. \textit{Contaminated checkpoints.} The models of \citet{oren2024proving} give randomised injection at $m\in\{1,2\}$; continued pretraining of Pythia-160M and 410M extends this to $m\le32$ across verbatim, evaluation-formatted, paraphrased \cite{yang2023rephrased} and source-document conditions, yielding a detection floor per format.

\textit{GSM8K against GSM1k.} GSM1k was commissioned to mirror GSM8K on human solve rate, solution steps and answer magnitude; its authors report accuracy drops of up to 8 points and a positive relationship, Spearman $r^2=0.36$, between a model's probability of generating a GSM8K example and its gap between the sets \cite{zhang2024gsm1k}. That evidence is behavioural and cannot separate item memorisation from stylistic overfitting or residual difficulty mismatch. RSCP supplies an independent representational line on the same split, and the four combinations of gap and contrast each license a different conclusion. One check costs nothing: computing $\BA_{\nuis}$ here audits whether the commissioned matching succeeded, which the benchmark's own checks do not establish.

\section{Significance}

The contribution is a protocol whose every component was adopted because a simpler alternative was measured and failed. Its central move is to stop assuming what a null depth profile looks like and measure it per model, using data the audit already requires: the placebo stops being a veto and becomes the baseline. The validation matters as much as the protocol, because three of the four discarded alternatives are what a careful reader would have tried first.

Two consequences reach further. The nuisance audit is a standalone result extending \citet{das2024blind} to splits they did not cover, needing forward passes only. And the most informative outcome may be negative: if the detection floor lies above the duplication counts that occur in practice, activation probing does not solve low-duplication contamination, and a negative result with a stated floor is worth more than another detector reporting high AUC on a split a blind classifier also solves.

Limits are stated plainly. RSCP is white-box only. Calibration is close but not exact, between \SlopeAAdjZero{} and \SlopeBAdjZero{} across the null cells. A weak surface key makes the instrument blind without announcing it. A null is weak evidence of absence, since linear probes lower-bound decodability. Calibration transfer across scales is assumed. And a developer wanting to hide contamination can delete a linearly encoded property without removing its behavioural benefit \cite{ravfogel2020inlp}, so this is an instrument for analysis, not an audit that survives an adversary. Above all: none of these measurements is evidence that a real transformer carries a familiarity direction.

\section{Follow-on Research}

The programme is pre-registered with falsification criteria. \textbf{H1} nuisance dominance on temporally split benchmarks; \textbf{H2} monotone exposure response fit by the calibration link; \textbf{H3} layer profile peaking in the middle-to-late third, not at the embedding layers; \textbf{H4} full-item extraction inflating the statistic relative to prefix extraction on sets differing in difficulty but not membership; \textbf{H5} decay with training steps elapsed since exposure on Pythia. Ahead of all of them sits \textbf{H0}, that the placebo contrast is null on real models; if surface decodability is not flat in depth, nothing downstream is interpretable.

Phase 0 is done and reported above. Five remain: the nuisance audit of existing splits, calibration by randomised injection, the public contaminated checkpoints, the GSM8K application with OLMo corpus search as partial ground truth, and robustness checks on probe family, extraction point, sample size and scale. Extraction and inference take minutes; continued pretraining dominates at an estimated 60 to 80 GPU-hours, roughly 100 in total, a fortnight on one consumer card. Released artefacts include cached activations, probe and nuisance weights, the bootstrap nulls rather than only the $p$-values, the injection corpora and checkpoints, and the Phase 0 implementation.

The decision rule is fixed in advance. If the exposure response holds and the floor is low, RSCP is a usable auditing instrument and the next step is deployment across benchmarks with commissioned twins. If the response holds but the floor is high, the floor itself is the contribution. If the response fails, the nuisance audit stands and the probe does not.

\renewcommand{\refname}{\normalsize\bfseries References}
\begin{multicols}{2}
\begin{thebibliography}{9}
\scriptsize
\setlength{\leftmargin}{0pt}
\setlength{\itemsep}{0pt}
\setlength{\parsep}{0pt}
\setlength{\labelsep}{3pt}

\bibitem[Biderman et al.(2023)]{biderman2023pythia} S.~Biderman et al. 2023. Pythia: A suite for analyzing LLMs across training and scaling. \emph{ICML}.

\bibitem[Braun(2026)]{braun2026crosslingual} F.~Braun. 2026. Interpretable cross-lingual alignment in small language models. % full title and venue

\bibitem[Das et al.(2024)]{das2024blind} D.~Das, J.~Zhang, F.~Tram\`er. 2024. Blind baselines beat membership inference attacks. \emph{arXiv:2406.16201}.

\bibitem[Duan et al.(2024)]{duan2024mia} M.~Duan et al. 2024. Do membership inference attacks work on large language models? \emph{COLM}.

\bibitem[Liu et al.(2024)]{liu2024probing} Z.~Liu et al. 2024. Probing language models for pre-training data detection. \emph{ACL}, 1576--1587.

\bibitem[Oren et al.(2024)]{oren2024proving} Y.~Oren et al. 2024. Proving test set contamination in black box LMs. \emph{ICLR}.

\bibitem[Ravfogel et al.(2020)]{ravfogel2020inlp} S.~Ravfogel et al. 2020. Null it out. \emph{ACL}.

\bibitem[Yang et al.(2023)]{yang2023rephrased} S.~Yang et al. 2023. Benchmark contamination with rephrased samples. \emph{arXiv:2311.04850}.

\bibitem[Zhang et al.(2024)]{zhang2024gsm1k} H.~Zhang et al. 2024. LLM performance on grade school arithmetic. \emph{NeurIPS D\&B}.

\end{thebibliography}
\end{multicols}

\end{document}
